\documentclass[12pt,a4paper]{article}

% =========================================================
% PAQUETES
% =========================================================
\usepackage[utf8]{inputenc}
\usepackage[T1]{fontenc}
\usepackage[spanish,es-nodecimaldot]{babel}
\usepackage{lmodern}
\usepackage{geometry}
\usepackage{setspace}
\usepackage{parskip}
\usepackage{enumitem}
\usepackage{longtable}
\usepackage{array}
\usepackage{booktabs}
\usepackage{tabularx}
\usepackage{xcolor}
\usepackage{titlesec}
\usepackage{fancyhdr}
\usepackage{hyperref}
\usepackage{graphicx}
\usepackage{float}
\usepackage{microtype}
\usepackage{amsmath}
\usepackage{amssymb}

% =========================================================
% CONFIGURACIÓN GENERAL
% =========================================================
\geometry{
top=2.5cm,
bottom=2.5cm,
left=2.8cm,
right=2.8cm
}

\setstretch{1.15}
\setlength{\parindent}{0pt}
\setlength{\parskip}{6pt}

\definecolor{primary}{HTML}{17365D}
\definecolor{secondary}{HTML}{2F75B5}
\definecolor{lightgray}{HTML}{F2F2F2}
\definecolor{darkgray}{HTML}{404040}
\definecolor{requirement}{HTML}{EAF2F8}
\definecolor{warning}{HTML}{FFF2CC}

\hypersetup{
colorlinks=true,
linkcolor=primary,
urlcolor=secondary,
citecolor=primary,
pdftitle={Especificación de Requisitos de Software},
pdfauthor={Proyecto de Asistencia Biométrica Escolar}
}

\titleformat{\section}
{\Large\bfseries\color{primary}}
{\thesection.}{0.5em}{}

\titleformat{\subsection}
{\large\bfseries\color{secondary}}
{\thesubsection.}{0.5em}{}

\titleformat{\subsubsection}
{\normalsize\bfseries\color{darkgray}}
{\thesubsubsection.}{0.5em}{}

\pagestyle{fancy}
\fancyhf{}
\fancyhead[L]{Plataforma de Asistencia Biométrica Escolar}
\fancyhead[R]{Especificación de Requisitos}
\fancyfoot[C]{\thepage}
\setlength{\headheight}{15pt}

% =========================================================
% COMANDOS PERSONALIZADOS
% =========================================================
\newcommand{\RF}[1]{\textbf{\color{primary}RF-#1}}
\newcommand{\RNF}[1]{\textbf{\color{secondary}RNF-#1}}
\newcommand{\RN}[1]{\textbf{\color{darkgray}RN-#1}}

\newenvironment{requisito}[1]
{
\begin{center}
\begin{minipage}{0.96\textwidth}
\colorbox{requirement}{
\begin{minipage}{0.94\textwidth}
\textbf{#1}
\end{minipage}
}
\vspace{2mm}
}
{
\end{minipage}
\end{center}
}

% =========================================================
% DOCUMENTO
% =========================================================
\begin{document}

% =========================================================
% PORTADA
% =========================================================
\begin{titlepage}
\centering

```
\vspace*{1.5cm}

{\Huge\bfseries\color{primary}
Especificación de Requisitos de Software\par}

\vspace{0.8cm}

{\LARGE
Plataforma de Asistencia Escolar Biométrica\\
y Gestión Académica\par}

\vspace{1.5cm}

\rule{\textwidth}{1.2pt}

\vspace{1cm}

\begin{tabular}{rl}
    \textbf{Tipo de documento:} & Especificación de Requisitos de Software\\
    \textbf{Versión:} & 1.0\\
    \textbf{Estado:} & Documento base para análisis y desarrollo\\
    \textbf{Plataformas objetivo:} & Android, iOS y dispositivos de acceso\\
    \textbf{Arquitectura propuesta:} & Aplicación móvil y servicios en la nube\\
\end{tabular}

\vfill

{\large
Documento estructurado con base en prácticas de especificación\\
de requisitos de software y criterios de calidad verificables.\par}

\vspace{1cm}

{\large \today\par}
```

\end{titlepage}

\tableofcontents
\newpage

% =========================================================
\section{Introducción}
% =========================================================

\subsection{Propósito}

El propósito de este documento es definir formalmente los requisitos funcionales,
no funcionales, reglas de negocio, actores, restricciones y criterios de calidad
de una plataforma de asistencia biométrica escolar y gestión académica.

El documento servirá como referencia para:

\begin{itemize}[leftmargin=*]
\item El análisis y diseño de la arquitectura de software.
\item La construcción del producto.
\item La planificación del desarrollo.
\item La creación del backlog del proyecto.
\item La elaboración de casos de uso.
\item El diseño de pruebas funcionales y no funcionales.
\item La validación del sistema por parte de la institución educativa.
\end{itemize}

\subsection{Alcance del sistema}

La plataforma permitirá administrar procesos académicos, asistencias,
calificaciones, evidencias, inscripciones, créditos, documentos escolares,
credenciales digitales y accesos físicos a la institución.

El sistema deberá funcionar principalmente mediante aplicaciones móviles para
Android e iOS. También deberá disponer de una modalidad de kiosco para los
dispositivos instalados en los accesos de la institución, tales como teléfonos,
tabletas, cámaras o equipos conectados a torniquetes.

Toda la información institucional deberá almacenarse en servicios de nube para
permitir su consulta desde diferentes ubicaciones y mantener la operación del
sistema las 24 horas del día.

\subsection{Objetivos generales}

\begin{itemize}[leftmargin=*]
\item Automatizar el registro de asistencia escolar.
\item Validar la identidad de los usuarios mediante mecanismos biométricos.
\item Reducir registros de asistencia fraudulentos.
\item Centralizar la información académica y administrativa.
\item Facilitar la gestión de alumnos, docentes y grupos.
\item Proporcionar trazabilidad sobre cambios administrativos sensibles.
\item Digitalizar las credenciales escolares y los permisos vehiculares.
\item Permitir el acceso institucional mediante mecanismos alternativos ante fallas.
\end{itemize}

\subsection{Definiciones y abreviaturas}

\begin{longtable}{p{3.5cm}p{10.5cm}}
\toprule
\textbf{Término} & \textbf{Definición}\
\midrule
RF & Requisito funcional. Define una operación o comportamiento que debe proporcionar el sistema.\
RNF & Requisito no funcional. Define una propiedad de calidad, restricción técnica o condición de operación.\
RN & Regla de negocio. Define una política o condición institucional que debe respetarse.\
API & Interfaz de programación utilizada para comunicar la aplicación móvil con los servicios del backend.\
JWT & Token firmado utilizado para autenticar y autorizar solicitudes hacia la API.\
CRUD & Operaciones de creación, consulta, actualización y eliminación de registros.\
Embedding biométrico & Representación matemática de las características faciales de una persona.\
Kiosco & Dispositivo o modalidad destinada exclusivamente al control de acceso físico.\
Control Escolar & Personal administrativo encargado de expedientes, inscripciones, pagos, calificaciones y créditos.\
Dirección & Rol administrativo con el mayor nivel de supervisión institucional.\
Fail-safe & Mecanismo alternativo utilizado para conservar una operación básica ante una falla.\
\bottomrule
\end{longtable}

% =========================================================
\section{Descripción general}
% =========================================================

\subsection{Perspectiva del producto}

El producto estará compuesto por los siguientes elementos:

\begin{enumerate}[leftmargin=*]
\item Aplicación móvil para alumnos.
\item Aplicación móvil para docentes.
\item Interfaz administrativa para Control Escolar.
\item Interfaz administrativa para Dirección.
\item Aplicación o modalidad de kiosco para el acceso físico.
\item API centralizada para la lógica de negocio.
\item Base de datos relacional alojada en la nube.
\item Servicio de almacenamiento de archivos.
\item Servicio de reconocimiento biométrico.
\item Servicio de notificaciones.
\item Servicio de recuperación de cuentas.
\end{enumerate}

La aplicación podrá utilizar un mismo proyecto multiplataforma y deberá
presentar diferentes módulos, menús y permisos dependiendo del rol autenticado.

\subsection{Actores del sistema}

Se consideran cinco tipos principales de usuario:

\begin{longtable}{p{3.5cm}p{10.5cm}}
\toprule
\textbf{Actor} & \textbf{Responsabilidad principal}\
\midrule
Alumno &
Consultar información académica, registrar asistencia, entregar evidencias,
realizar trámites, utilizar su credencial digital y consultar sus créditos.\

Maestro &
Administrar grupos, actividades, asistencias y calificaciones de las materias
que tenga asignadas.\

Control Escolar &
Gestionar alumnos, docentes, grupos, inscripciones, pagos, créditos,
documentación y correcciones administrativas autorizadas.\

Dirección &
Supervisar integralmente la institución y ejecutar las operaciones de Control
Escolar con privilegios adicionales de seguimiento y análisis.\

Kiosco de acceso &
Registrar y validar el acceso físico de los usuarios mediante reconocimiento
facial, códigos QR, PIN u otros mecanismos autorizados.\
\bottomrule
\end{longtable}

\subsection{Restricciones generales}

\begin{itemize}[leftmargin=*]
\item La aplicación deberá ejecutarse en Android e iOS.
\item La información institucional deberá permanecer centralizada en la nube.
\item La aplicación no deberá depender exclusivamente del almacenamiento local.
\item Las operaciones deberán estar restringidas por roles y permisos.
\item Los datos biométricos deberán tratarse como información sensible.
\item El sistema deberá mantener trazabilidad sobre modificaciones administrativas.
\item La operación estará disponible permanentemente, salvo suspensiones autorizadas.
\end{itemize}

\subsection{Suposiciones y dependencias}

\begin{itemize}[leftmargin=*]
\item La institución proporcionará los datos oficiales de alumnos, docentes,
carreras, grupos, salones y periodos escolares.

```
\item Los dispositivos móviles utilizados para asistencia biométrica deberán
contar con sensores compatibles o mecanismos seguros de autenticación local.

\item La geolocalización dependerá de los servicios GPS y permisos del dispositivo.

\item El reconocimiento facial dependerá de la calidad de las cámaras, iluminación,
conectividad y capacidad de procesamiento disponible.

\item Los servicios de SMS, correo, almacenamiento y notificaciones podrán depender
de proveedores externos.

\item La institución definirá los calendarios, fechas límite, periodos de evaluación
y políticas de acceso.
```

\end{itemize}

% =========================================================
\section{Requisitos funcionales}
% =========================================================

\subsection{Módulo de autenticación y control de acceso}

\begin{requisito}{\RF{1.1.1} --- Inicio de sesión unificado}
El sistema deberá proporcionar una única pantalla de inicio de sesión para todos
los usuarios.

Después de validar las credenciales, el sistema deberá identificar el rol del
usuario y mostrar únicamente la interfaz, funcionalidades y datos autorizados
para dicho rol.

\textbf{Actores:} Alumno, Maestro, Control Escolar, Dirección y Kiosco.

\textbf{Prioridad:} Alta.

\textbf{Criterio de aceptación:} Un usuario autenticado no podrá visualizar ni
ejecutar funciones correspondientes a otro rol sin autorización explícita.
\end{requisito}

\begin{requisito}{\RF{1.1.2} --- Identificación de usuarios}
El sistema deberá aceptar los siguientes identificadores iniciales:

\begin{itemize}[leftmargin=*]
\item Alumno: matrícula escolar.
\item Maestro: número de empleado o número de maestro.
\item Control Escolar: cuenta administrativa asignada por la institución.
\item Dirección: nombre completo o cuenta administrativa asignada.
\item Kiosco: identificador institucional único del dispositivo o ubicación.
\end{itemize}

Los identificadores deberán ser únicos dentro del dominio correspondiente.
\end{requisito}

\begin{requisito}{\RF{1.1.3} --- Generación de credenciales temporales}
El sistema deberá permitir la creación de contraseñas temporales conforme a las
reglas iniciales proporcionadas por la institución:

\begin{itemize}[leftmargin=*]
\item Dirección: contraseña temporal \texttt{admin}.
\item Control Escolar: nombre de la escuela seguido de \texttt{1#}.
\item Maestro: últimos dos dígitos del número de empleado, siglas de la carrera
y el carácter \texttt{#}.
\item Alumno: contraseña temporal \texttt{123456789}.
\item Kiosco: nombre de la escuela seguido de \texttt{12}.
\end{itemize}

Por razones de seguridad, todas las contraseñas anteriores deberán considerarse
\textbf{contraseñas temporales de aprovisionamiento} y no contraseñas permanentes.

El sistema deberá exigir su cambio durante el primer inicio de sesión.
\end{requisito}

\begin{requisito}{\RF{1.1.4} --- Cambio obligatorio de contraseña}
El sistema deberá obligar a los usuarios personales a cambiar su contraseña
temporal durante el primer inicio de sesión.

La nueva contraseña deberá cumplir la política de seguridad institucional.

Hasta completar el cambio, el usuario no deberá tener acceso a las demás
funciones de la plataforma.
\end{requisito}

\begin{requisito}{\RF{1.1.5} --- Registro de medios de recuperación}
Durante el primer inicio de sesión, el alumno y el maestro deberán registrar al
menos uno de los siguientes medios de recuperación:

\begin{itemize}[leftmargin=*]
\item Dirección de correo electrónico válida.
\item Número de teléfono celular válido.
\end{itemize}

El sistema deberá verificar el medio registrado mediante un código de un solo uso.
\end{requisito}

\begin{requisito}{\RF{1.1.6} --- Recuperación de contraseña}
El sistema deberá permitir a alumnos y maestros recuperar su cuenta mediante un
código de verificación enviado al correo electrónico o teléfono previamente
registrado.

El código deberá:

\begin{itemize}[leftmargin=*]
\item Tener una vigencia limitada.
\item Ser de un solo uso.
\item Invalidarse después de varios intentos fallidos.
\item No mostrarse en registros visibles para usuarios administrativos.
\end{itemize}
\end{requisito}

\begin{requisito}{\RF{1.1.7} --- Gestión de sesiones}
El sistema deberá permitir:

\begin{itemize}[leftmargin=*]
\item Iniciar sesión.
\item Cerrar sesión.
\item Revocar sesiones activas.
\item Invalidar sesiones cuando una cuenta sea suspendida.
\item Renovar una sesión mediante mecanismos seguros.
\end{itemize}
\end{requisito}

\begin{requisito}{\RF{1.1.8} --- Bloqueo por intentos fallidos}
El sistema deberá bloquear temporalmente una cuenta cuando se detecte un número
configurable de intentos de autenticación fallidos.

El sistema deberá registrar el evento y permitir el desbloqueo automático o
administrativo conforme a las políticas institucionales.
\end{requisito}

% ---------------------------------------------------------
\subsection{Módulo del alumno}
% ---------------------------------------------------------

\begin{requisito}{\RF{1.2.1} --- Registro inicial de fotografía}
Durante el primer inicio de sesión, el alumno deberá capturar una fotografía
facial desde la aplicación.

La aplicación deberá comprobar, como mínimo:

\begin{itemize}[leftmargin=*]
\item Presencia de un rostro.
\item Calidad mínima de iluminación.
\item Nitidez suficiente.
\item Ausencia de múltiples rostros.
\item Correspondencia con una captura en tiempo real.
\end{itemize}

La captura deberá enviarse al servicio biométrico para generar la plantilla o
embedding utilizado en el reconocimiento facial.
\end{requisito}

\begin{requisito}{\RF{1.2.2} --- Perfil del alumno}
El alumno deberá poder consultar su perfil institucional, incluyendo:

\begin{itemize}[leftmargin=*]
\item Nombre completo.
\item Matrícula.
\item Fotografía.
\item Carrera.
\item Semestre.
\item Grupo.
\item Estatus académico.
\item Correo y teléfono vinculados.
\end{itemize}

La modificación de información institucional sensible deberá requerir
autorización de Control Escolar.
\end{requisito}

\begin{requisito}{\RF{1.2.3} --- Dashboard académico}
El sistema deberá mostrar al alumno un panel académico con:

\begin{itemize}[leftmargin=*]
\item Materias del periodo actual.
\item Calificaciones parciales.
\item Calificaciones finales.
\item Promedio del periodo.
\item Promedio general.
\item Alertas académicas.
\item Actividades próximas a vencer.
\item Resumen de asistencias.
\end{itemize}
\end{requisito}

\begin{requisito}{\RF{1.2.4} --- Consulta de historial académico}
El alumno deberá poder consultar su historial académico por:

\begin{itemize}[leftmargin=*]
\item Semestre.
\item Periodo escolar.
\item Materia.
\item Calificación.
\item Estatus de acreditación.
\item Créditos obtenidos.
\end{itemize}

El historial deberá mostrar únicamente datos oficiales y validados.
\end{requisito}

\begin{requisito}{\RF{1.2.5} --- Registro de asistencia en aula}
El alumno deberá poder registrar su asistencia desde la aplicación utilizando
la autenticación biométrica disponible en su dispositivo.

El registro deberá considerar:

\begin{itemize}[leftmargin=*]
\item Identidad del alumno autenticado.
\item Materia y grupo.
\item Fecha y hora.
\item Ubicación aproximada.
\item Ventana de tiempo autorizada.
\item Sesión de clase activa.
\item Resultado de la validación biométrica local.
\end{itemize}

La aplicación no deberá extraer ni almacenar directamente la huella digital o
información interna de Face ID. Solamente deberá utilizar el resultado seguro de
autenticación proporcionado por Android o iOS.
\end{requisito}

\begin{requisito}{\RF{1.2.6} --- Validación por geolocalización}
Antes de aceptar una asistencia, el sistema deberá verificar que el dispositivo
se encuentre dentro de una zona geográfica autorizada para la institución o el
salón correspondiente.

La tolerancia geográfica deberá ser configurable por la institución.
\end{requisito}

\begin{requisito}{\RF{1.2.7} --- Prevención de asistencias duplicadas}
El sistema no deberá permitir que un alumno registre más de una asistencia para
la misma sesión de clase.

Los intentos posteriores deberán registrarse como eventos de auditoría sin
duplicar la asistencia.
\end{requisito}

\begin{requisito}{\RF{1.2.8} --- Consulta de asistencias}
El alumno deberá poder consultar sus asistencias clasificadas como:

\begin{itemize}[leftmargin=*]
\item Presente.
\item Retardo.
\item Falta.
\item Justificada.
\item Pendiente de validación.
\item Modificada por el docente.
\end{itemize}
\end{requisito}

\begin{requisito}{\RF{1.2.9} --- Entrega de evidencias}
El alumno deberá poder cargar evidencias asociadas a actividades académicas,
incluyendo documentos PDF e imágenes.

Por cada entrega, el sistema deberá almacenar:

\begin{itemize}[leftmargin=*]
\item Alumno.
\item Actividad.
\item Archivo o archivos entregados.
\item Fecha y hora exactas.
\item Estado de entrega.
\item Indicador de entrega extemporánea.
\item Calificación.
\item Retroalimentación del docente.
\end{itemize}
\end{requisito}

\begin{requisito}{\RF{1.2.10} --- Sustitución de evidencias}
El alumno podrá sustituir una evidencia únicamente cuando:

\begin{itemize}[leftmargin=*]
\item La fecha límite no haya vencido.
\item El docente haya habilitado nuevos intentos.
\item La actividad permita múltiples entregas.
\end{itemize}

El sistema deberá conservar el historial de versiones conforme a la política de
retención institucional.
\end{requisito}

\begin{requisito}{\RF{1.2.11} --- Consulta de actividades}
El alumno deberá poder consultar las actividades asignadas por sus maestros,
incluyendo:

\begin{itemize}[leftmargin=*]
\item Nombre.
\item Descripción.
\item Tipo de actividad.
\item Materia.
\item Valor porcentual o puntaje.
\item Fecha de publicación.
\item Fecha límite.
\item Archivos adjuntos.
\item Estado de entrega.
\end{itemize}
\end{requisito}

\begin{requisito}{\RF{1.2.12} --- Credencial digital}
El sistema deberá generar una credencial digital equivalente a la credencial
física institucional.

La credencial deberá incluir:

\begin{itemize}[leftmargin=*]
\item Fotografía.
\item Nombre completo.
\item Matrícula.
\item Carrera.
\item Periodo de vigencia.
\item Código QR o código de barras.
\item Identificador verificable.
\item Estatus de vigencia.
\end{itemize}

El código mostrado deberá poder validarse contra el servidor para detectar
credenciales vencidas, revocadas o alteradas.
\end{requisito}

\begin{requisito}{\RF{1.2.13} --- Consulta de créditos}
El alumno deberá poder consultar:

\begin{itemize}[leftmargin=*]
\item Créditos académicos obtenidos.
\item Créditos académicos pendientes.
\item Créditos complementarios obtenidos.
\item Créditos complementarios pendientes.
\item Servicio Social.
\item Prácticas Profesionales.
\item Actividades complementarias.
\item Certificaciones de idioma.
\end{itemize}
\end{requisito}

\begin{requisito}{\RF{1.2.14} --- Mecanismo alternativo de acceso}
La aplicación deberá generar un mecanismo temporal de acceso cuando el sistema
principal de reconocimiento de entrada no esté disponible.

El mecanismo podrá consistir en:

\begin{itemize}[leftmargin=*]
\item Código QR dinámico.
\item PIN temporal.
\item Código firmado de un solo uso.
\end{itemize}

El código deberá tener vigencia limitada y no deberá poder reutilizarse después
de ser validado.
\end{requisito}

\begin{requisito}{\RF{1.2.15} --- Solicitud de documentos escolares}
El alumno deberá poder solicitar documentos institucionales, tales como:

\begin{itemize}[leftmargin=*]
\item Constancia de estudios.
\item Historial académico.
\item Constancia de inscripción.
\item Documentos adicionales configurados por la institución.
\end{itemize}

El sistema deberá mostrar el estado de la solicitud y notificar al alumno cuando
el documento esté disponible.
\end{requisito}

\begin{requisito}{\RF{1.2.16} --- Reinscripción y selección de materias}
El alumno deberá poder iniciar un proceso de reinscripción y seleccionar las
materias disponibles conforme a:

\begin{itemize}[leftmargin=*]
\item Carrera.
\item Semestre.
\item Plan de estudios.
\item Materias acreditadas.
\item Prerrequisitos.
\item Cupo disponible.
\item Compatibilidad de horarios.
\end{itemize}
\end{requisito}

\begin{requisito}{\RF{1.2.17} --- Carga de comprobante de pago}
Durante la reinscripción, el alumno deberá poder cargar una imagen o documento
correspondiente al comprobante de pago.

El registro deberá permanecer con estatus \texttt{Pendiente de validación} hasta
que Control Escolar confirme el pago.
\end{requisito}

\begin{requisito}{\RF{1.2.18} --- Confirmación presencial de inscripción}
La inscripción no deberá considerarse definitiva hasta que Control Escolar
confirme la recepción y validación de los documentos físicos requeridos.

Los estados mínimos del proceso serán:

\begin{itemize}[leftmargin=*]
\item Borrador.
\item Pago pendiente.
\item Pago en revisión.
\item Documentos físicos pendientes.
\item Inscripción confirmada.
\item Rechazada.
\end{itemize}
\end{requisito}

% ---------------------------------------------------------
\subsection{Módulo del maestro}
% ---------------------------------------------------------

\begin{requisito}{\RF{1.3.1} --- Gestión de perfil docente}
El maestro deberá poder consultar su información institucional y modificar su
fotografía de perfil.

Los cambios en nombre, número de empleado, carrera o adscripción deberán ser
autorizados por Control Escolar.
\end{requisito}

\begin{requisito}{\RF{1.3.2} --- Consulta de carga académica}
El maestro deberá poder consultar:

\begin{itemize}[leftmargin=*]
\item Materias asignadas.
\item Grupos asignados.
\item Salones.
\item Horarios.
\item Periodos escolares.
\item Listas de alumnos.
\end{itemize}
\end{requisito}

\begin{requisito}{\RF{1.3.3} --- Calendario docente}
El sistema deberá mostrar al maestro un calendario con:

\begin{itemize}[leftmargin=*]
\item Días de clase.
\item Horarios.
\item Actividades.
\item Evaluaciones.
\item Fechas límite de captura.
\item Periodos vacacionales.
\item Días de suspensión.
\end{itemize}
\end{requisito}

\begin{requisito}{\RF{1.3.4} --- Creación de actividades}
El maestro deberá poder crear actividades para sus grupos, incluyendo:

\begin{itemize}[leftmargin=*]
\item Actividades en clase.
\item Tareas.
\item Proyectos.
\item Exámenes.
\item Prácticas.
\item Participaciones.
\item Otros tipos configurables.
\end{itemize}

Cada actividad deberá incluir nombre, descripción, materia, grupo, valor,
fecha de publicación y fecha límite.
\end{requisito}

\begin{requisito}{\RF{1.3.5} --- Configuración de rubros de evaluación}
El maestro deberá poder definir los rubros de evaluación de cada materia y
parcial.

La suma de los porcentajes deberá ser igual al 100% antes de cerrar la
configuración del periodo.

Ejemplo:

\begin{itemize}[leftmargin=*]
\item Examen: 40%.
\item Tareas: 30%.
\item Proyecto: 30%.
\end{itemize}
\end{requisito}

\begin{requisito}{\RF{1.3.6} --- Captura de calificaciones por actividad}
El maestro deberá poder registrar una calificación y retroalimentación para cada
entrega realizada por un alumno.

El sistema deberá validar que la calificación se encuentre dentro de la escala
permitida por la institución.
\end{requisito}

\begin{requisito}{\RF{1.3.7} --- Cálculo de calificación parcial}
El sistema deberá calcular automáticamente la calificación parcial con base en:

\begin{itemize}[leftmargin=*]
\item Calificaciones de actividades.
\item Valor de cada actividad.
\item Rubros de evaluación.
\item Reglas de redondeo.
\item Actividades anuladas o justificadas.
\end{itemize}

El maestro deberá revisar y confirmar el resultado antes de su publicación.
\end{requisito}

\begin{requisito}{\RF{1.3.8} --- Captura de calificaciones finales}
El maestro deberá poder registrar y enviar calificaciones finales durante el
periodo autorizado por Control Escolar.

Una vez vencido el periodo, la modificación deberá quedar bloqueada para el
maestro.
\end{requisito}

\begin{requisito}{\RF{1.3.9} --- Gestión de asistencias}
El maestro deberá poder consultar en tiempo real las asistencias registradas por
los alumnos en cada sesión de clase.

El panel deberá mostrar:

\begin{itemize}[leftmargin=*]
\item Alumno.
\item Hora del registro.
\item Resultado de validación.
\item Ubicación validada.
\item Estatus de asistencia.
\item Observaciones.
\end{itemize}
\end{requisito}

\begin{requisito}{\RF{1.3.10} --- Validación manual de asistencia}
El maestro podrá validar o modificar la asistencia de un alumno cuando exista
una causa justificada.

La modificación deberá registrar:

\begin{itemize}[leftmargin=*]
\item Maestro responsable.
\item Fecha y hora.
\item Estado anterior.
\item Estado nuevo.
\item Motivo.
\end{itemize}
\end{requisito}

\begin{requisito}{\RF{1.3.11} --- Recordatorios de fechas límite}
El sistema deberá enviar recordatorios al maestro antes del vencimiento de las
fechas de captura de calificaciones parciales y finales.

La anticipación de los recordatorios deberá ser configurable.
\end{requisito}

\begin{requisito}{\RF{1.3.12} --- Consulta de entregas}
El maestro deberá poder consultar, descargar y revisar las evidencias entregadas
por los alumnos de sus grupos.

El maestro no deberá poder consultar archivos de grupos o materias que no tenga
asignados.
\end{requisito}

% ---------------------------------------------------------
\subsection{Módulo de Control Escolar}
% ---------------------------------------------------------

\begin{requisito}{\RF{1.4.1} --- Gestión de alumnos}
Control Escolar deberá poder crear, consultar, actualizar, suspender y reactivar
registros de alumnos.

La eliminación física de expedientes deberá estar restringida. Preferentemente,
el sistema deberá utilizar baja lógica para conservar la trazabilidad.
\end{requisito}

\begin{requisito}{\RF{1.4.2} --- Gestión de maestros}
Control Escolar deberá poder registrar y administrar:

\begin{itemize}[leftmargin=*]
\item Datos del maestro.
\item Número de empleado.
\item Carrera o departamento.
\item Materias asignadas.
\item Grupos.
\item Horarios.
\item Estatus laboral dentro del sistema.
\end{itemize}
\end{requisito}

\begin{requisito}{\RF{1.4.3} --- Importación masiva}
Control Escolar deberá poder importar alumnos y maestros mediante archivos
estructurados, como CSV o Excel.

Antes de guardar la información, el sistema deberá:

\begin{itemize}[leftmargin=*]
\item Validar columnas obligatorias.
\item Detectar registros duplicados.
\item Identificar datos inválidos.
\item Mostrar un resumen de errores.
\item Solicitar confirmación.
\end{itemize}
\end{requisito}

\begin{requisito}{\RF{1.4.4} --- Gestión de grupos y asignaciones}
Control Escolar deberá poder:

\begin{itemize}[leftmargin=*]
\item Crear grupos.
\item Asignar alumnos.
\item Asignar maestros.
\item Asignar materias.
\item Asignar salones.
\item Configurar horarios.
\item Definir periodos escolares.
\end{itemize}
\end{requisito}

\begin{requisito}{\RF{1.4.5} --- Validación de pagos}
Control Escolar deberá poder consultar los comprobantes cargados por los alumnos
y clasificarlos como:

\begin{itemize}[leftmargin=*]
\item Pendiente.
\item Aprobado.
\item Rechazado.
\item Requiere aclaración.
\end{itemize}

Cuando se rechace un comprobante, deberá registrarse el motivo.
\end{requisito}

\begin{requisito}{\RF{1.4.6} --- Confirmación de inscripción}
Control Escolar deberá poder confirmar la inscripción cuando:

\begin{itemize}[leftmargin=*]
\item El pago esté validado.
\item Los documentos físicos hayan sido entregados.
\item La selección de materias sea válida.
\item No existan bloqueos administrativos.
\end{itemize}
\end{requisito}

\begin{requisito}{\RF{1.4.7} --- Corrección extemporánea de calificaciones}
Control Escolar deberá poder modificar una calificación después de la fecha
límite únicamente cuando cuente con el permiso correspondiente.

La operación deberá exigir:

\begin{itemize}[leftmargin=*]
\item Selección del alumno.
\item Materia.
\item Periodo.
\item Calificación anterior.
\item Calificación nueva.
\item Motivo obligatorio.
\item Confirmación de la operación.
\end{itemize}
\end{requisito}

\begin{requisito}{\RF{1.4.8} --- Asignación de créditos complementarios}
Control Escolar deberá poder registrar créditos y avances relacionados con:

\begin{itemize}[leftmargin=*]
\item Servicio Social.
\item Prácticas Profesionales.
\item Idioma inglés.
\item Actividades culturales.
\item Actividades deportivas.
\item Certificaciones.
\item Otros créditos definidos por la institución.
\end{itemize}
\end{requisito}

\begin{requisito}{\RF{1.4.9} --- Gestión de documentos escolares}
Control Escolar deberá poder revisar, aprobar, rechazar y emitir las solicitudes
de documentos realizadas por los alumnos.

Cada cambio de estado deberá notificarse al solicitante.
\end{requisito}

\begin{requisito}{\RF{1.4.10} --- Gestión de calendarios y periodos}
Control Escolar deberá poder configurar:

\begin{itemize}[leftmargin=*]
\item Periodos escolares.
\item Fechas de inicio y fin.
\item Periodos de evaluación.
\item Fechas límite de calificaciones.
\item Fechas de reinscripción.
\item Días festivos.
\item Suspensiones.
\end{itemize}
\end{requisito}

\begin{requisito}{\RF{1.4.11} --- Suspensión de cuentas}
Control Escolar deberá poder suspender temporalmente una cuenta por causas
administrativas.

La suspensión deberá incluir:

\begin{itemize}[leftmargin=*]
\item Motivo.
\item Fecha de inicio.
\item Fecha de finalización opcional.
\item Usuario responsable.
\item Alcance de la suspensión.
\end{itemize}
\end{requisito}

% ---------------------------------------------------------
\subsection{Módulo de Dirección}
% ---------------------------------------------------------

\begin{requisito}{\RF{1.5.1} --- Funciones administrativas heredadas}
Dirección deberá disponer de las funciones asignadas a Control Escolar, sujetas
a una política de autorización de mayor nivel.

Todas las operaciones deberán identificar claramente al usuario que las ejecutó.
\end{requisito}

\begin{requisito}{\RF{1.5.2} --- Seguimiento integral de alumnos}
Dirección deberá poder consultar el expediente integral de un alumno,
incluyendo:

\begin{itemize}[leftmargin=*]
\item Historial académico.
\item Asistencias.
\item Calificaciones.
\item Créditos.
\item Trámites.
\item Incidencias.
\item Estatus de inscripción.
\item Riesgos académicos.
\end{itemize}
\end{requisito}

\begin{requisito}{\RF{1.5.3} --- Seguimiento de maestros}
Dirección deberá poder consultar indicadores e incidencias relacionadas con los
maestros, tales como:

\begin{itemize}[leftmargin=*]
\item Grupos asignados.
\item Cumplimiento de captura de calificaciones.
\item Registro de asistencias.
\item Actividades creadas.
\item Incidencias administrativas.
\item Observaciones institucionales.
\end{itemize}
\end{requisito}

\begin{requisito}{\RF{1.5.4} --- Seguimiento de grupos}
Dirección deberá poder analizar el desempeño de cada grupo mediante indicadores
como:

\begin{itemize}[leftmargin=*]
\item Promedio general.
\item Índice de aprobación.
\item Índice de reprobación.
\item Asistencia promedio.
\item Alumnos en riesgo.
\item Materias con mayor reprobación.
\end{itemize}
\end{requisito}

\begin{requisito}{\RF{1.5.5} --- Panel de indicadores}
El sistema deberá proporcionar a Dirección un panel con indicadores académicos,
administrativos y de acceso físico.

Los indicadores deberán poder filtrarse por periodo, carrera, semestre, grupo,
materia y maestro.
\end{requisito}

\begin{requisito}{\RF{1.5.6} --- Generación de reportes}
Dirección deberá poder generar reportes en formatos configurables, incluyendo
PDF y hojas de cálculo.

Los reportes podrán incluir:

\begin{itemize}[leftmargin=*]
\item Asistencias.
\item Calificaciones.
\item Rendimiento académico.
\item Alumnos en riesgo.
\item Incidencias.
\item Accesos físicos.
\item Créditos.
\item Inscripciones.
\end{itemize}
\end{requisito}

\begin{requisito}{\RF{1.5.7} --- Suspensión general del servicio}
Dirección deberá poder activar una suspensión programada de acceso a la
plataforma durante periodos vacacionales, días festivos o contingencias.

La suspensión deberá permitir definir su alcance:

\begin{itemize}[leftmargin=*]
\item Solo alumnos.
\item Alumnos y maestros.
\item Acceso físico.
\item Operaciones académicas específicas.
\item Plataforma completa.
\end{itemize}

Las funciones administrativas críticas y de emergencia deberán permanecer
disponibles para los usuarios autorizados.
\end{requisito}

% ---------------------------------------------------------
\subsection{Módulo de kiosco y acceso físico}
% ---------------------------------------------------------

\begin{requisito}{\RF{1.6.1} --- Autenticación del dispositivo de acceso}
Cada dispositivo utilizado en la entrada deberá autenticarse mediante una cuenta
o credencial técnica única.

No deberá utilizarse una misma credencial compartida para todos los dispositivos
en producción.

La cuenta deberá estar asociada con:

\begin{itemize}[leftmargin=*]
\item Institución.
\item Puerta o acceso.
\item Dispositivo.
\item Estado.
\item Fecha de activación.
\end{itemize}
\end{requisito}

\begin{requisito}{\RF{1.6.2} --- Escaneo facial continuo}
El kiosco deberá mantener activa la cámara durante su horario de operación y
detectar rostros presentes frente al dispositivo.

El sistema deberá evitar enviar solicitudes repetidas del mismo rostro dentro de
un intervalo configurable.
\end{requisito}

\begin{requisito}{\RF{1.6.3} --- Verificación de presencia real}
El sistema de acceso deberá utilizar mecanismos de detección de vida o
\textit{liveness detection} para reducir intentos de acceso mediante:

\begin{itemize}[leftmargin=*]
\item Fotografías impresas.
\item Fotografías mostradas en pantallas.
\item Videos pregrabados.
\item Máscaras u otros medios de suplantación.
\end{itemize}
\end{requisito}

\begin{requisito}{\RF{1.6.4} --- Reconocimiento de identidad}
Cuando se detecte un rostro válido, el kiosco deberá enviar al servicio
biométrico la información necesaria para compararlo con las plantillas
registradas.

El sistema deberá devolver:

\begin{itemize}[leftmargin=*]
\item Identidad encontrada o no encontrada.
\item Nivel de confianza.
\item Estatus del usuario.
\item Vigencia de la credencial.
\item Decisión de acceso.
\end{itemize}
\end{requisito}

\begin{requisito}{\RF{1.6.5} --- Decisión de acceso}
Cuando la identidad sea válida y el usuario tenga autorización, el kiosco deberá
mostrar:

\begin{center}
\textbf{ACCESO PERMITIDO}
\end{center}

La pantalla podrá mostrar temporalmente:

\begin{itemize}[leftmargin=*]
\item Fotografía.
\item Nombre.
\item Matrícula o número de empleado.
\item Tipo de usuario.
\end{itemize}

Cuando la identidad no sea válida, deberá mostrar:

\begin{center}
\textbf{ACCESO DENEGADO}
\end{center}
\end{requisito}

\begin{requisito}{\RF{1.6.6} --- Activación del mecanismo físico}
Cuando el acceso sea autorizado, el sistema deberá enviar la señal correspondiente
al torniquete, puerta o mecanismo de apertura integrado.

La apertura deberá tener una duración limitada y deberá registrarse en la bitácora.
\end{requisito}

\begin{requisito}{\RF{1.6.7} --- Registro de entradas y salidas}
El sistema deberá registrar cada intento de acceso con:

\begin{itemize}[leftmargin=*]
\item Usuario identificado, cuando corresponda.
\item Fecha y hora.
\item Puerta o torniquete.
\item Dispositivo.
\item Tipo de evento: entrada o salida.
\item Método utilizado.
\item Resultado.
\item Motivo de rechazo, cuando corresponda.
\end{itemize}
\end{requisito}

\begin{requisito}{\RF{1.6.8} --- Validación de QR o PIN alternativo}
El kiosco deberá poder validar los códigos QR dinámicos o PIN temporales
generados por la aplicación.

El código deberá rechazarse cuando:

\begin{itemize}[leftmargin=*]
\item Haya expirado.
\item Ya haya sido utilizado.
\item No corresponda al usuario.
\item Haya sido revocado.
\item No esté firmado correctamente.
\end{itemize}
\end{requisito}

\begin{requisito}{\RF{1.6.9} --- Operación ante pérdida de conectividad}
Cuando exista una pérdida temporal de conectividad, el dispositivo deberá mostrar
claramente que se encuentra en modo degradado.

La institución podrá habilitar una lista cifrada y limitada de credenciales
temporales para continuar operaciones básicas.

Todos los eventos generados sin conexión deberán sincronizarse cuando se
restablezca el servicio.
\end{requisito}

\begin{requisito}{\RF{1.6.10} --- Restricción funcional del kiosco}
El usuario de kiosco solamente deberá acceder a las funciones necesarias para:

\begin{itemize}[leftmargin=*]
\item Detectar y validar rostros.
\item Validar códigos alternativos.
\item Mostrar el resultado de acceso.
\item Registrar entradas y salidas.
\item Consultar el estado de conectividad.
\item Ejecutar diagnósticos autorizados.
\end{itemize}

El kiosco no deberá poder consultar calificaciones, expedientes, pagos,
documentos ni información académica adicional.
\end{requisito}

% ---------------------------------------------------------
\subsection{Módulo de credenciales y control vehicular}
% ---------------------------------------------------------

\begin{requisito}{\RF{1.7.1} --- Habilitación manual de pase vehicular}
Control Escolar deberá poder habilitar manualmente un pase vehicular para un
alumno, maestro o usuario administrativo.

El pase no deberá generarse automáticamente.
\end{requisito}

\begin{requisito}{\RF{1.7.2} --- Registro de vehículo}
El sistema deberá permitir registrar:

\begin{itemize}[leftmargin=*]
\item Tipo de vehículo.
\item Marca.
\item Modelo.
\item Color.
\item Placas.
\item Nombre del propietario o responsable.
\item Matrícula o número de empleado.
\item Fotografía del usuario.
\item Fecha de inicio de vigencia.
\item Fecha de vencimiento.
\end{itemize}
\end{requisito}

\begin{requisito}{\RF{1.7.3} --- Tipos de vehículo}
El sistema deberá admitir, como mínimo, los siguientes tipos:

\begin{itemize}[leftmargin=*]
\item Motocicleta.
\item Automóvil.
\item Camioneta.
\end{itemize}

La institución podrá configurar tipos adicionales.
\end{requisito}

\begin{requisito}{\RF{1.7.4} --- Credencial vehicular digital}
Cuando el pase se encuentre autorizado, la aplicación deberá mostrar una
credencial vehicular digital con:

\begin{itemize}[leftmargin=*]
\item Fotografía del usuario.
\item Nombre.
\item Matrícula o número de empleado.
\item Placas.
\item Tipo de vehículo.
\item Vigencia.
\item Código QR verificable.
\item Estatus del permiso.
\end{itemize}
\end{requisito}

\begin{requisito}{\RF{1.7.5} --- Revocación de pase vehicular}
Control Escolar y Dirección deberán poder suspender o revocar un permiso
vehicular.

La revocación deberá surtir efecto inmediatamente en los dispositivos conectados.
\end{requisito}

\begin{requisito}{\RF{1.7.6} --- Registro de acceso vehicular}
El sistema deberá registrar la entrada y salida de vehículos, incluyendo:

\begin{itemize}[leftmargin=*]
\item Usuario.
\item Vehículo.
\item Placas.
\item Fecha y hora de entrada.
\item Fecha y hora de salida.
\item Punto de acceso.
\item Resultado de la validación.
\end{itemize}
\end{requisito}

% ---------------------------------------------------------
\subsection{Módulo de notificaciones}
% ---------------------------------------------------------

\begin{requisito}{\RF{1.8.1} --- Notificaciones internas}
El sistema deberá proporcionar un centro de notificaciones dentro de la
aplicación.

Las notificaciones deberán indicar su estado como leída o no leída.
\end{requisito}

\begin{requisito}{\RF{1.8.2} --- Notificaciones push}
El sistema deberá enviar notificaciones push para eventos relevantes, como:

\begin{itemize}[leftmargin=*]
\item Actividad asignada.
\item Fecha de entrega próxima.
\item Calificación publicada.
\item Cambio de asistencia.
\item Documento disponible.
\item Pago rechazado.
\item Inscripción confirmada.
\item Fecha límite de captura docente.
\end{itemize}
\end{requisito}

\begin{requisito}{\RF{1.8.3} --- Comunicados institucionales}
Control Escolar y Dirección deberán poder enviar comunicados a:

\begin{itemize}[leftmargin=*]
\item Toda la institución.
\item Una carrera.
\item Un semestre.
\item Un grupo.
\item Maestros.
\item Alumnos.
\item Usuarios específicos.
\end{itemize}
\end{requisito}

% =========================================================
\section{Reglas de negocio}
% =========================================================

\begin{requisito}{\RN{01} --- Contraseñas iniciales}
Las contraseñas generadas durante el alta de usuarios serán temporales y deberán
cambiarse durante el primer inicio de sesión.
\end{requisito}

\begin{requisito}{\RN{02} --- Restricción de calificaciones}
Los maestros solamente podrán modificar calificaciones dentro del periodo
habilitado por Control Escolar.
\end{requisito}

\begin{requisito}{\RN{03} --- Modificación extemporánea}
Toda modificación de calificaciones fuera del periodo deberá requerir un motivo
y generar un registro de auditoría.
\end{requisito}

\begin{requisito}{\RN{04} --- Confirmación de inscripción}
La carga de un comprobante de pago no confirmará automáticamente una inscripción.
\end{requisito}

\begin{requisito}{\RN{05} --- Documentación física}
La inscripción se confirmará únicamente después de que Control Escolar valide
los documentos físicos requeridos.
\end{requisito}

\begin{requisito}{\RN{06} --- Asistencia válida}
Una asistencia será válida cuando se encuentre dentro del horario, ubicación y
sesión de clase autorizados.
\end{requisito}

\begin{requisito}{\RN{07} --- Asistencia única}
Un alumno no podrá registrar más de una asistencia para la misma sesión.
\end{requisito}

\begin{requisito}{\RN{08} --- Credencial digital}
Una credencial digital vencida, suspendida o revocada no permitirá el acceso.
\end{requisito}

\begin{requisito}{\RN{09} --- Pase vehicular}
El pase vehicular solamente será válido durante el periodo autorizado por
Control Escolar.
\end{requisito}

\begin{requisito}{\RN{10} --- Privilegios de Dirección}
Dirección podrá ejecutar funciones de Control Escolar y consultar información
institucional consolidada.
\end{requisito}

\begin{requisito}{\RN{11} --- Acceso mínimo del kiosco}
El kiosco tendrá permisos mínimos y no podrá acceder a información académica.
\end{requisito}

\begin{requisito}{\RN{12} --- Suspensión programada}
Las suspensiones generales deberán registrar responsable, motivo, fecha de inicio,
fecha de finalización y alcance.
\end{requisito}

% =========================================================
\section{Requisitos no funcionales}
% =========================================================

\subsection{Arquitectura y tecnología}

\begin{requisito}{\RNF{2.1.1} --- Aplicación multiplataforma}
La aplicación móvil deberá ser compatible con Android e iOS.

Se propone utilizar React Native con Expo, siempre que las funciones biométricas,
de cámara y geolocalización requeridas sean compatibles con la estrategia de
compilación seleccionada.
\end{requisito}

\begin{requisito}{\RNF{2.1.2} --- Backend}
El backend deberá implementar una API desacoplada de la aplicación cliente.

Se propone utilizar Python y FastAPI por su compatibilidad con operaciones
asíncronas, servicios biométricos y documentación automática de API.
\end{requisito}

\begin{requisito}{\RNF{2.1.3} --- Base de datos}
La información transaccional deberá almacenarse en una base de datos relacional.

Se propone PostgreSQL por sus capacidades de integridad, concurrencia,
transacciones, extensibilidad y auditoría.
\end{requisito}

\begin{requisito}{\RNF{2.1.4} --- Arquitectura modular}
El backend deberá separar, al menos, los siguientes dominios:

\begin{itemize}[leftmargin=*]
\item Identidad y acceso.
\item Gestión académica.
\item Asistencias.
\item Evaluaciones.
\item Archivos.
\item Biometría.
\item Inscripciones.
\item Credenciales.
\item Control vehicular.
\item Notificaciones.
\item Auditoría.
\end{itemize}
\end{requisito}

\begin{requisito}{\RNF{2.1.5} --- Almacenamiento de archivos}
Las fotografías, comprobantes, evidencias y documentos no deberán almacenarse
directamente como archivos binarios dentro de las tablas transaccionales.

Deberán almacenarse en un servicio de objetos, como:

\begin{itemize}[leftmargin=*]
\item Amazon S3.
\item Cloudinary.
\item Firebase Storage.
\item Servicio compatible con S3.
\end{itemize}

La base de datos almacenará la referencia, metadatos, propietario, tipo,
integridad y permisos del archivo.
\end{requisito}

\begin{requisito}{\RNF{2.1.6} --- Servicios en la nube}
El backend, la base de datos y el almacenamiento deberán desplegarse en
infraestructura de nube.

No deberán depender de un único equipo local de la institución para operar.
\end{requisito}

\begin{requisito}{\RNF{2.1.7} --- Separación de ambientes}
El proyecto deberá contar, como mínimo, con ambientes separados para:

\begin{itemize}[leftmargin=*]
\item Desarrollo.
\item Pruebas.
\item Producción.
\end{itemize}

Los datos reales de producción no deberán utilizarse directamente en desarrollo.
\end{requisito}

\subsection{Disponibilidad y continuidad}

\begin{requisito}{\RNF{2.2.1} --- Disponibilidad}
El sistema deberá estar disponible las 24 horas del día, los 7 días de la semana,
excepto durante:

\begin{itemize}[leftmargin=*]
\item Mantenimientos programados.
\item Suspensiones institucionales autorizadas.
\item Eventos de fuerza mayor.
\end{itemize}
\end{requisito}

\begin{requisito}{\RNF{2.2.2} --- Nivel objetivo de disponibilidad}
La disponibilidad mensual objetivo del backend deberá ser igual o superior al
99.5%, excluyendo mantenimientos programados previamente comunicados.
\end{requisito}

\begin{requisito}{\RNF{2.2.3} --- Respaldo de base de datos}
La base de datos deberá contar con respaldos automáticos.

Como mínimo, deberán implementarse:

\begin{itemize}[leftmargin=*]
\item Respaldo diario.
\item Retención configurable.
\item Cifrado de respaldos.
\item Pruebas periódicas de restauración.
\end{itemize}
\end{requisito}

\begin{requisito}{\RNF{2.2.4} --- Recuperación ante desastres}
La arquitectura deberá definir:

\begin{itemize}[leftmargin=*]
\item Objetivo de tiempo de recuperación, RTO.
\item Objetivo de punto de recuperación, RPO.
\item Procedimiento de restauración.
\item Responsables.
\item Pruebas periódicas.
\end{itemize}

Como valor inicial, se propone un RTO máximo de cuatro horas y un RPO máximo de
veinticuatro horas, sujetos a aprobación institucional.
\end{requisito}

\subsection{Seguridad}

\begin{requisito}{\RNF{2.3.1} --- Control de acceso basado en roles}
El sistema deberá implementar control de acceso basado en roles y permisos.

Cada endpoint de la API deberá verificar tanto la autenticación como la
autorización del usuario.
\end{requisito}

\begin{requisito}{\RNF{2.3.2} --- Autenticación mediante tokens}
La comunicación autenticada entre las aplicaciones y la API deberá utilizar
tokens firmados con expiración limitada.

Los tokens de renovación deberán almacenarse y transmitirse mediante mecanismos
seguros.
\end{requisito}

\begin{requisito}{\RNF{2.3.3} --- Protección de contraseñas}
Las contraseñas no deberán almacenarse en texto plano.

Deberán protegerse mediante un algoritmo de derivación de claves adecuado, como
Argon2id o bcrypt, con parámetros de costo configurables.
\end{requisito}

\begin{requisito}{\RNF{2.3.4} --- Política de contraseñas}
Las contraseñas permanentes deberán cumplir, como mínimo:

\begin{itemize}[leftmargin=*]
\item Longitud mínima configurable.
\item Combinación de diferentes tipos de caracteres.
\item Rechazo de contraseñas comunes.
\item Prohibición de reutilizar contraseñas recientes.
\item Cambio obligatorio cuando exista una recuperación de cuenta.
\end{itemize}
\end{requisito}

\begin{requisito}{\RNF{2.3.5} --- Cifrado en tránsito}
Toda comunicación entre clientes, API, base de datos y almacenamiento deberá
protegerse mediante TLS.

No se permitirán credenciales ni datos biométricos transmitidos por conexiones
sin cifrado.
\end{requisito}

\begin{requisito}{\RNF{2.3.6} --- Cifrado en reposo}
Los datos sensibles, respaldos y objetos almacenados deberán permanecer cifrados
en reposo mediante las capacidades del proveedor de infraestructura.
\end{requisito}

\begin{requisito}{\RNF{2.3.7} --- Datos biométricos}
El sistema deberá minimizar el almacenamiento de imágenes faciales.

Para el reconocimiento, deberá almacenar preferentemente plantillas biométricas
o embeddings cifrados, junto con los metadatos indispensables.

La fotografía institucional utilizada para la credencial deberá tratarse
separadamente de la plantilla biométrica.
\end{requisito}

\begin{requisito}{\RNF{2.3.8} --- Consentimiento biométrico}
Antes del registro biométrico, el sistema deberá presentar un aviso de privacidad
y obtener el consentimiento correspondiente conforme a las políticas y
obligaciones legales aplicables.
\end{requisito}

\begin{requisito}{\RNF{2.3.9} --- Prevención de suplantación}
El reconocimiento facial deberá incluir controles de detección de vida y
protección contra ataques de presentación.

El sistema deberá registrar intentos sospechosos sin exponer información
biométrica al usuario final.
\end{requisito}

\begin{requisito}{\RNF{2.3.10} --- Protección de archivos}
Los archivos privados no deberán exponerse mediante direcciones públicas
permanentes.

El acceso deberá realizarse mediante permisos, URLs firmadas y tiempos de
expiración.
\end{requisito}

\begin{requisito}{\RNF{2.3.11} --- Validación de archivos}
Antes de aceptar un archivo, el sistema deberá validar:

\begin{itemize}[leftmargin=*]
\item Tipo MIME.
\item Extensión.
\item Tamaño.
\item Integridad.
\item Asociación con el usuario.
\item Posibles contenidos maliciosos.
\end{itemize}
\end{requisito}

\begin{requisito}{\RNF{2.3.12} --- Gestión de secretos}
Las claves de API, contraseñas de base de datos, certificados y secretos no
deberán almacenarse directamente en el código fuente.

Deberán administrarse mediante variables de entorno o un servicio especializado
de gestión de secretos.
\end{requisito}

\begin{requisito}{\RNF{2.3.13} --- Limitación de solicitudes}
La API deberá aplicar límites de solicitudes para reducir ataques de fuerza
bruta, abuso de endpoints y saturación del sistema.
\end{requisito}

\subsection{Auditoría y trazabilidad}

\begin{requisito}{\RNF{2.4.1} --- Bitácora de auditoría}
El sistema deberá mantener una bitácora de operaciones sensibles.

La bitácora deberá registrar:

\begin{itemize}[leftmargin=*]
\item Usuario responsable.
\item Acción.
\item Recurso afectado.
\item Fecha y hora.
\item Dirección IP.
\item Dispositivo o cliente.
\item Valor anterior.
\item Valor nuevo.
\item Motivo, cuando sea obligatorio.
\end{itemize}
\end{requisito}

\begin{requisito}{\RNF{2.4.2} --- Auditoría de calificaciones}
Toda modificación extemporánea de calificaciones deberá generar automáticamente
un registro de auditoría desde la base de datos o desde una capa transaccional
protegida.

La aplicación no deberá permitir al usuario administrativo eliminar o alterar
dicho registro.
\end{requisito}

\begin{requisito}{\RNF{2.4.3} --- Integridad de auditoría}
Los registros de auditoría deberán ser de solo inserción para los usuarios de la
aplicación.

Las correcciones sobre una bitácora deberán realizarse mediante nuevos eventos,
sin sobrescribir los anteriores.
\end{requisito}

\begin{requisito}{\RNF{2.4.4} --- Sincronización horaria}
Todos los servidores y dispositivos de acceso deberán mantener sus relojes
sincronizados con una fuente horaria confiable.

Las fechas deberán almacenarse en UTC y presentarse de acuerdo con la zona
horaria institucional.
\end{requisito}

\subsection{Rendimiento}

\begin{requisito}{\RNF{2.5.1} --- Tiempo de reconocimiento facial}
Bajo condiciones normales de conectividad y carga, el sistema de acceso deberá
emitir una decisión en un tiempo máximo de 2.5 segundos desde la captura válida
del rostro.
\end{requisito}

\begin{requisito}{\RNF{2.5.2} --- Tiempo de respuesta de consultas}
Las operaciones comunes de consulta deberán responder, en el percentil 95, en
menos de dos segundos, sin considerar la descarga de archivos pesados.
\end{requisito}

\begin{requisito}{\RNF{2.5.3} --- Operaciones pesadas}
La carga y procesamiento de archivos pesados deberá ejecutarse de manera
asíncrona y no deberá bloquear la navegación principal de la aplicación.
\end{requisito}

\begin{requisito}{\RNF{2.5.4} --- Procesamiento en segundo plano}
Las siguientes operaciones podrán ejecutarse mediante colas de trabajo:

\begin{itemize}[leftmargin=*]
\item Procesamiento de fotografías.
\item Generación de embeddings.
\item Envío de correos.
\item Envío de notificaciones.
\item Generación de reportes.
\item Procesamiento de archivos.
\end{itemize}
\end{requisito}

\begin{requisito}{\RNF{2.5.5} --- Concurrencia}
La infraestructura deberá dimensionarse con base en:

\begin{itemize}[leftmargin=*]
\item Número total de alumnos.
\item Número de maestros.
\item Accesos simultáneos.
\item Cantidad de torniquetes.
\item Horarios de máxima entrada.
\item Periodos de publicación de calificaciones.
\end{itemize}

Las pruebas de carga deberán utilizar escenarios reales proporcionados por la
institución.
\end{requisito}

\subsection{Usabilidad y accesibilidad}

\begin{requisito}{\RNF{2.6.1} --- Diseño responsivo}
La interfaz deberá adaptarse a diferentes tamaños de pantalla, incluyendo
teléfonos pequeños, teléfonos de gama media y tabletas.
\end{requisito}

\begin{requisito}{\RNF{2.6.2} --- Legibilidad}
El sistema deberá utilizar tipografías, tamaños, espacios y contrastes que
permitan una lectura clara.

Las pantallas de acceso deberán mantener visibilidad suficiente bajo iluminación
exterior.
\end{requisito}

\begin{requisito}{\RNF{2.6.3} --- Retroalimentación inmediata}
La aplicación deberá informar claramente al usuario cuando una operación:

\begin{itemize}[leftmargin=*]
\item Está en proceso.
\item Fue completada.
\item Falló.
\item Requiere corrección.
\item Está pendiente de validación.
\end{itemize}
\end{requisito}

\begin{requisito}{\RNF{2.6.4} --- Mensajes de error}
Los mensajes de error deberán ser comprensibles y no deberán revelar detalles
internos, consultas SQL, rutas del servidor, tokens ni información sensible.
\end{requisito}

\begin{requisito}{\RNF{2.6.5} --- Accesibilidad}
La interfaz deberá considerar:

\begin{itemize}[leftmargin=*]
\item Compatibilidad con escalado de texto.
\item Etiquetas para lectores de pantalla.
\item Controles táctiles con tamaño suficiente.
\item Contraste adecuado.
\item Navegación consistente.
\item Información no dependiente exclusivamente del color.
\end{itemize}
\end{requisito}

\subsection{Compatibilidad}

\begin{requisito}{\RNF{2.7.1} --- Versiones de Android}
La aplicación deberá establecer una versión mínima de Android conforme a las
capacidades biométricas y de seguridad requeridas.

La matriz de compatibilidad deberá definirse antes de iniciar las pruebas de
aceptación.
\end{requisito}

\begin{requisito}{\RNF{2.7.2} --- Versiones de iOS}
La aplicación deberá establecer una versión mínima de iOS conforme a las
capacidades de Face ID, Touch ID, cámara, notificaciones y almacenamiento seguro.
\end{requisito}

\begin{requisito}{\RNF{2.7.3} --- Dispositivos de kiosco}
El módulo de kiosco deberá poder ejecutarse en dispositivos autorizados con
cámara compatible y conexión segura.

Cuando se integre con torniquetes, deberá definirse un protocolo de comunicación
documentado.
\end{requisito}

\subsection{Mantenibilidad y calidad}

\begin{requisito}{\RNF{2.8.1} --- Código modular}
El código deberá organizarse en módulos con responsabilidades claramente
definidas y bajo acoplamiento.
\end{requisito}

\begin{requisito}{\RNF{2.8.2} --- Documentación de API}
La API deberá contar con documentación actualizada de:

\begin{itemize}[leftmargin=*]
\item Endpoints.
\item Métodos.
\item Parámetros.
\item Respuestas.
\item Códigos de error.
\item Permisos requeridos.
\end{itemize}
\end{requisito}

\begin{requisito}{\RNF{2.8.3} --- Pruebas automatizadas}
El proyecto deberá contar con pruebas automatizadas para:

\begin{itemize}[leftmargin=*]
\item Autenticación.
\item Autorización.
\item Cálculo de calificaciones.
\item Registro de asistencia.
\item Inscripción.
\item Auditoría.
\item Validación de credenciales.
\end{itemize}
\end{requisito}

\begin{requisito}{\RNF{2.8.4} --- Integración continua}
El repositorio deberá utilizar un proceso de integración continua que ejecute
validaciones de código y pruebas antes de desplegar una nueva versión.
\end{requisito}

\begin{requisito}{\RNF{2.8.5} --- Observabilidad}
El sistema deberá disponer de:

\begin{itemize}[leftmargin=*]
\item Registros estructurados.
\item Métricas.
\item Monitoreo de disponibilidad.
\item Alertas.
\item Seguimiento de errores.
\item Identificadores de correlación.
\end{itemize}
\end{requisito}

\subsection{Privacidad y retención}

\begin{requisito}{\RNF{2.9.1} --- Minimización de datos}
El sistema deberá recopilar únicamente la información necesaria para las
funciones académicas, administrativas y de seguridad autorizadas.
\end{requisito}

\begin{requisito}{\RNF{2.9.2} --- Retención de información}
La institución deberá definir periodos de retención para:

\begin{itemize}[leftmargin=*]
\item Expedientes académicos.
\item Evidencias.
\item Fotografías.
\item Plantillas biométricas.
\item Registros de acceso.
\item Auditorías.
\item Comprobantes de pago.
\end{itemize}
\end{requisito}

\begin{requisito}{\RNF{2.9.3} --- Baja de datos biométricos}
Cuando un usuario deje de pertenecer a la institución, el sistema deberá permitir
revocar su acceso y eliminar o anonimizar sus datos biométricos conforme a la
política legal de retención aplicable.
\end{requisito}

% =========================================================
\section{Matriz de permisos}
% =========================================================

\renewcommand{\arraystretch}{1.3}

\begin{longtable}{
p{5.2cm}
>{\centering\arraybackslash}p{1.4cm}
>{\centering\arraybackslash}p{1.4cm}
>{\centering\arraybackslash}p{1.8cm}
>{\centering\arraybackslash}p{1.5cm}
>{\centering\arraybackslash}p{1.4cm}
}
\toprule
\textbf{Función} &
\textbf{Alumno} &
\textbf{Maestro} &
\textbf{Control Escolar} &
\textbf{Dirección} &
\textbf{Kiosco}\
\midrule
\endfirsthead

\toprule
\textbf{Función} &
\textbf{Alumno} &
\textbf{Maestro} &
\textbf{Control Escolar} &
\textbf{Dirección} &
\textbf{Kiosco}\
\midrule
\endhead

Consultar calificaciones propias & Sí & No & Sí & Sí & No\
Registrar asistencia propia & Sí & No & No & No & No\
Validar asistencia de un grupo & No & Sí & Sí & Sí & No\
Crear actividades & No & Sí & No & No & No\
Calificar actividades & No & Sí & No & No & No\
Modificar calificaciones en periodo & No & Sí & Supervisión & Supervisión & No\
Modificar calificaciones extemporáneas & No & No & Sí & Sí & No\
Gestionar alumnos & No & No & Sí & Sí & No\
Gestionar maestros & No & No & Sí & Sí & No\
Validar pagos & No & No & Sí & Sí & No\
Confirmar inscripciones & No & No & Sí & Sí & No\
Asignar créditos & No & No & Sí & Sí & No\
Consultar indicadores globales & No & No & Limitado & Sí & No\
Autorizar pase vehicular & No & No & Sí & Sí & No\
Mostrar credencial vehicular propia & Sí & Sí & Sí & Sí & No\
Validar acceso físico & No & No & No & No & Sí\
Suspender acceso general & No & No & Según permiso & Sí & No\
Consultar auditorías & No & No & Limitado & Sí & No\

\bottomrule
\end{longtable}

% =========================================================
\section{Estados principales del sistema}
% =========================================================

\subsection{Estados de una cuenta}

\begin{itemize}[leftmargin=*]
\item Pendiente de activación.
\item Contraseña temporal.
\item Activa.
\item Bloqueada temporalmente.
\item Suspendida.
\item Inactiva.
\item Revocada.
\end{itemize}

\subsection{Estados de una inscripción}

\begin{itemize}[leftmargin=*]
\item Borrador.
\item Pago pendiente.
\item Pago en revisión.
\item Pago rechazado.
\item Documentos físicos pendientes.
\item En validación.
\item Confirmada.
\item Rechazada.
\item Cancelada.
\end{itemize}

\subsection{Estados de una entrega}

\begin{itemize}[leftmargin=*]
\item No entregada.
\item Entregada a tiempo.
\item Entregada fuera de tiempo.
\item En revisión.
\item Calificada.
\item Requiere corrección.
\item Anulada.
\end{itemize}

\subsection{Estados de una asistencia}

\begin{itemize}[leftmargin=*]
\item Pendiente.
\item Presente.
\item Retardo.
\item Falta.
\item Justificada.
\item Rechazada por ubicación.
\item Rechazada por horario.
\item Modificada manualmente.
\end{itemize}

\subsection{Estados de una credencial vehicular}

\begin{itemize}[leftmargin=*]
\item Pendiente.
\item Activa.
\item Suspendida.
\item Vencida.
\item Revocada.
\end{itemize}

% =========================================================
\section{Criterios generales de aceptación}
% =========================================================

El sistema podrá considerarse funcionalmente aceptable cuando se compruebe que:

\begin{enumerate}[leftmargin=*]
\item Cada usuario accede únicamente a las funciones autorizadas por su rol.

```
\item Las contraseñas temporales deben cambiarse durante el primer acceso.

\item Los alumnos pueden consultar su información académica y entregar
evidencias.

\item Los maestros pueden gestionar actividades, asistencias y calificaciones
de sus grupos.

\item Control Escolar puede administrar alumnos, pagos, inscripciones,
documentos y créditos.

\item Dirección puede supervisar la información institucional y generar
reportes consolidados.

\item El kiosco puede identificar usuarios y emitir una decisión de acceso.

\item Cada modificación sensible genera un registro de auditoría.

\item Los archivos se almacenan fuera de la base de datos transaccional.

\item Las comunicaciones están cifradas.

\item Las contraseñas no se almacenan en texto plano.

\item Los datos biométricos se encuentran protegidos y sujetos a políticas de
privacidad.

\item El acceso alternativo mediante QR o PIN funciona durante las
contingencias autorizadas.

\item Los permisos vehiculares pueden ser emitidos, consultados y revocados.

\item El sistema satisface los tiempos de respuesta establecidos para las
operaciones críticas.
```

\end{enumerate}

% =========================================================
\section{Riesgos arquitectónicos identificados}
% =========================================================

\begin{longtable}{p{4.5cm}p{5cm}p{4.5cm}}
\toprule
\textbf{Riesgo} & \textbf{Impacto} & \textbf{Tratamiento recomendado}\
\midrule

Contraseñas iniciales predecibles &
Acceso no autorizado a cuentas nuevas &
Forzar cambio inmediato, limitar vigencia y notificar la activación.\

Uso de una sola cuenta para todos los kioscos &
Imposibilidad de identificar el dispositivo comprometido &
Asignar una credencial técnica diferente a cada equipo.\

Falsificación de ubicación &
Registro fraudulento de asistencia &
Combinar geolocalización, sesión de clase, biometría y controles antifraude.\

Uso de fotografías para engañar al reconocimiento &
Acceso físico no autorizado &
Implementar detección de vida y umbrales de confianza.\

Pérdida de conectividad &
Interrupción del acceso a la institución &
Implementar modo degradado controlado y sincronización posterior.\

Almacenamiento excesivo de fotografías &
Riesgo de privacidad y aumento de costos &
Minimizar imágenes, utilizar plantillas biométricas y políticas de retención.\

Carga masiva incorrecta &
Corrupción o duplicación de datos &
Validar previamente, generar reporte de errores y usar transacciones.\

Modificaciones administrativas sin trazabilidad &
Pérdida de integridad académica &
Utilizar auditoría inmutable y motivos obligatorios.\

Sobrecarga durante horarios de entrada &
Filas y denegaciones por tiempo de espera &
Realizar pruebas de carga y escalar horizontalmente el servicio biométrico.\

Dependencia de un proveedor de nube &
Dificultad de migración o interrupción del servicio &
Utilizar estándares abiertos, respaldos exportables y almacenamiento compatible.\

\bottomrule
\end{longtable}

% =========================================================
\section{Priorización inicial}
% =========================================================

\subsection{Prioridad crítica}

\begin{itemize}[leftmargin=*]
\item Autenticación y autorización.
\item Gestión de alumnos y maestros.
\item Registro y validación de asistencia.
\item Reconocimiento facial de acceso.
\item Gestión de materias, grupos y horarios.
\item Captura de calificaciones.
\item Auditoría de modificaciones.
\item Infraestructura en la nube.
\end{itemize}

\subsection{Prioridad alta}

\begin{itemize}[leftmargin=*]
\item Evidencias académicas.
\item Recuperación de cuentas.
\item Credencial digital.
\item Inscripciones y pagos.
\item Créditos complementarios.
\item Notificaciones.
\item Códigos alternativos de acceso.
\end{itemize}

\subsection{Prioridad media}

\begin{itemize}[leftmargin=*]
\item Reportes avanzados.
\item Indicadores de riesgo académico.
\item Control vehicular.
\item Analítica institucional.
\item Exportaciones personalizadas.
\end{itemize}

% =========================================================
\section{Recomendación de arquitectura lógica}
% =========================================================

Se recomienda iniciar con un \textbf{monolito modular} desplegado en la nube,
en lugar de implementar microservicios desde la primera versión.

El monolito modular permitirá:

\begin{itemize}[leftmargin=*]
\item Reducir complejidad operativa.
\item Mantener transacciones consistentes.
\item Facilitar las pruebas.
\item Acelerar el desarrollo inicial.
\item Separar los dominios mediante módulos internos.
\item Extraer servicios independientes en el futuro cuando exista una necesidad
comprobada.
\end{itemize}

El servicio biométrico y los procesos en segundo plano podrán separarse desde el
inicio debido a sus necesidades específicas de cómputo y escalabilidad.

Una distribución lógica inicial sería:

\begin{enumerate}[leftmargin=*]
\item Aplicación móvil React Native.
\item Interfaz de kiosco.
\item API FastAPI.
\item Módulo de autenticación.
\item Módulo académico.
\item Módulo de asistencias.
\item Módulo de inscripciones.
\item Módulo de credenciales.
\item Servicio biométrico.
\item Cola de trabajos asíncronos.
\item PostgreSQL.
\item Almacenamiento de objetos.
\item Servicio de monitoreo y auditoría.
\end{enumerate}

% =========================================================
\section{Conclusión}
% =========================================================

La plataforma propuesta centraliza los procesos de control académico, asistencia
biométrica, evaluación, inscripción, credenciales y acceso físico de la
institución.

Los requisitos establecidos en este documento deberán utilizarse como base para:

\begin{itemize}[leftmargin=*]
\item Elaborar diagramas de casos de uso.
\item Diseñar el modelo de datos.
\item Definir la arquitectura de componentes.
\item Construir historias de usuario.
\item Estimar el proyecto.
\item Diseñar prototipos de interfaz.
\item Preparar pruebas de aceptación.
\item Definir un plan de implementación por fases.
\end{itemize}

Cualquier modificación posterior deberá registrarse mediante control de
versiones y mantener la trazabilidad entre el requisito, su diseño, su
implementación y sus pruebas.

\end{document}
